\begin{solucion}
Definamos primeramente los números triangulares
\[
  T_n=\frac{n(n+1)}{2},\qquad n\in\mathbb{N},
\]
y la función
\[
  C(n,k)\;:=\;(n-2k)\,\binom{n}{k},\qquad 0\le k\le n.
\]
Nuestro objetivo es probar que, para todo $n\ge 1$, se cumple la identidad
\begin{equation}\label{eq:id}
  \sum_{k=0}^{n-1}
        C\bigl(T_n,k\bigr)\,C\bigl(T_{n+1},k\bigr)
  \;=\;
  \frac{n^{3}-2n^{2}+4n}{n+2}\,
  \binom{T_n}{n}\,\binom{T_{n+1}}{n}.
\end{equation}
A continuación se reproduce la demostración \emph{computacional} utilizando los algoritmos de
\textsc{Gosper} y \textsc{Zeilberger} en \textit{Mathematica}. El cuaderno ejecutado se adjunta como
\href{run:Taller7_8.nb}{\texttt{Taller7\_8.nb}}.

% =====================================================================
\textbf{\textit{Mathematica}}

\begin{lstlisting}[caption={Definición de los objetos básicos}]
(* Números triangulares *)
T[n_] := n (n + 1)/2;

(* Función C(n,k) = (n - 2 k) Binomial[n, k] *)
c[n_, k_] := (n - 2 k) Binomial[n, k];

(* Sumando de interés *)
F[n_, k_] := c[T[n], k] * c[T[n + 1], k];
\end{lstlisting}

% ---------------------------------------------------------------------
\textbf{Aplicación del algoritmo de \textsc{Gosper}}

Usamos el algoritmo de \textsc{Gosper} para verificar que $F(n,k)$ es hipergeométrico en $k$.

\begin{lstlisting}[caption={Algoritmo de Gosper}]
Gosper[F[n, k], k]
{1/4 (-2+4 k-3 n-n^2) (4 k-n-n^2) Binomial[1/2 n (1+n),k] Binomial[1/2 (1+n) (2+n),k]==Subscript[\[CapitalDelta], k][-((k^2 (4-4 k+2 n+n^2) Binomial[1/2 n (1+n),k] Binomial[1/2 (1+n) (2+n),k])/(-2 n-n^2))]}
(* ==> {Gk, Δk}  (certificado telescopico) *)
\end{lstlisting}

Lo cual da como resultado la relación:
\begin{align*}
    \frac{1}{4} \left(4 k-n^2-3 n-2\right) \left(4 k-n^2-n\right) \binom{\frac{1}{2} n (n+1)}{k} \binom{\frac{1}{2} (n+1) (n+2)}{k} \\=\Delta _k\left(-\frac{k^2 \left(-4 k+n^2+2 n+4\right) \binom{\frac{1}{2} n (n+1)}{k} \binom{\frac{1}{2} (n+1) (n+2)}{k}}{-n^2-2 n}\right)
\end{align*}

Por lo tanto, concluimos que $F(n,k)$ se puede expresar como una anti-diferencia:
\[
  F(n,k)\;=\;\Delta_k G(n,k),
\]
donde
\[
  G(n,k)\;:=\;
  -\frac{k^{2}\bigl(4-4k+2n+n^{2}\bigr)}
         {-2n-n^{2}}\,
   \binom{\tfrac12 n(1+n)}{k}\,
   \binom{\tfrac12 (1+n)(2+n)}{k}.
\]

% ---------------------------------------------------------------------
\textbf{Certificado telescópico $G(n,k)$}

La llamada \verb|show[R] // FullSimplify| devuelve la función racional

\begin{lstlisting}[caption={Certificado }]
show[R] // FullSimplify
(4 k^2 (-4+4 k-n (2+n)))/((-4 k+n+n^2) (-4 k+(1+n) (2+n)))
\end{lstlisting}
Dando como resultado el certificado:
\[
  R(n,k)=\frac{4k^{2}\bigl(-4+4k-n(2+n)\bigr)}
               {(-4k+n+n^{2})(-4k+(1+n)(2+n))}.
\]
A partir de $R$ definimos el certificado
\[
  G(n,k)\;:=\;
  -\frac{k^{2}\bigl(4-4k+2n+n^{2}\bigr)}
         {-2n-n^{2}}\,
   \binom{\tfrac12 n(1+n)}{k}\,
   \binom{\tfrac12 (1+n)(2+n)}{k}.
\]

% ---------------------------------------------------------------------
\textbf{Verificación de la telescopica}

Comprobamos que
\[
  F(n,k)\;=\;G(n,k+1)-G(n,k),
\]
lo cual implica que la suma deseada se colapsa a los valores de $G$ en los extremos:
\[
  S(n)=G(n,n)-G(n,0).
\]
En \textit{Mathematica}:

\begin{lstlisting}[caption={Chequeo de telescopía}]
FullSimplify[F[n, k] - (G[n, k + 1] - G[n, k])]
(* ==> 0 *)
\end{lstlisting}

% ---------------------------------------------------------------------
\textbf{Evaluación de los bordes}

\begin{align*}
  G(n,n) &= \frac{n\,(4-2n+n^{2})}{n+2}\,
            \binom{\tfrac12 n(1+n)}{n}\,
            \binom{\tfrac12 (1+n)(2+n)}{n}, \\[6pt]
  G(n,0) &= 0.
\end{align*}
Por tanto
\[
  S(n)=G(n,n)-G(n,0)=
  \frac{n^{3}-2n^{2}+4n}{n+2}\,
  \binom{T_n}{n}\,\binom{T_{n+1}}{n},
\]
obteniendo exactamente el lado derecho de \eqref{eq:id}. La igualdad queda demostrada.

\end{solucion}